\section{A sufficient completion theorem and the remaining programme}
\label{sec:completion}

The purpose of a completion criterion is to identify precisely what the
remaining estimates have to deliver. It is not a theorem that assumes
Yang--Mills existence and then renames it. Its useful content is the
operator-theoretic assembly: once the scalar vacuum representation and
common-domain estimates are available, adding fields does not require
a new independent Euclidean reconstruction at every engineering degree.
The actual construction of those fields and their Yang--Mills
identification remain substantive premises.

\subsection{Completion in one vacuum representation}

Let \(V_D\) be increasing finite-dimensional label spaces for a chosen
basis of gauge-invariant \(SU(2)\) curvature/jet polynomials of engineering
dimension at most \(D\). Fix the algebraic relations, adjoints,
derivatives, and allowed renormalization mixing before defining these
spaces. The scalar family is among the labels. Bare multiplication of
symbols is not assumed to equal multiplication of renormalized fields at
the same point.

\begin{theorem}[Common-space completion criterion]\label{thm:completion}
Suppose a scalar vacuum theory \((\HH,U,\Omega)\) has been constructed,
where \(U\) is strongly continuous and Poincar\'e covariant, its spectrum
lies in the forward cone, and \(\Omega\) is its unique invariant unit
vector. Let \(H\geq0\) be the Hamiltonian and put
\(\DD_\infty=\bigcap_{r\geq0}\Dom(1+H)^r\).
For every \(D\), suppose that the actual lattice/fusion construction
gives linear operator distributions \(\Phi_D(v,f)\), \(v\in V_D\),
compatible under increasing \(D\), with the following properties.
\begin{enumerate}[label=(\alph*)]
 \item The fields and their adjoint companions preserve \(\DD_\infty\).
 For each \(r\geq0\) there exist \(p,s,C\), uniform on the finite label
 space, for which
 \begin{equation}
 \|(1+H)^r\Phi_D(v,f)(1+H)^{-p}\|
       \leq C|v|p_s(f). \label{eq:energybound}
 \end{equation}
 Here \(p_s\) is a Schwartz seminorm. The fields are jointly local on
 \(\DD_\infty\), and the scalar subfamily makes \(\Omega\) cyclic.
 \item The fields intertwine \(U\) with the declared tensor
 representations on their label spaces and obey the declared adjoint
 and derivative relations.
 \item The field map is nondegenerate on the chosen operator quotient
 and satisfies the claimed renormalized local identities, including
 the equation-of-motion/contact and Ward identities in their declared
 schemes. Its ultraviolet correlations and local operator products,
 with the fixed coupling and mixing convention, have the Yang--Mills
 coefficients and controlled remainders at each requested finite
 perturbative order. The stress-tensor Ward normalization is included
 when claiming that part of the identification.
 \item Some centered field creates a vacuum vector with a nonzero
 component outside the translation-invariant subspace.
\end{enumerate}
Then these data define a nontrivial local \(SU(2)\) observable quantum
field theory with the stated Yang--Mills identification. No mass gap
is implied.
\end{theorem}
\begin{proof}
The vacuum is in \(\DD_\infty\). Because that domain is invariant,
every finite product of smeared fields on the vacuum is defined.
Apply \eqref{eq:energybound} successively, choosing the input exponent
for one factor to absorb the output exponent required by the next.
For every fixed finite list of labels this gives
\begin{equation}
 |\langle\Omega,\Phi(v_1,f_1)\cdots\Phi(v_n,f_n)\Omega\rangle|
 \leq C_{D,n}\prod_{j=1}^n p_{s_j}(f_j). \label{eq:wordbound}
\end{equation}
The Schwartz kernel theorem yields a tempered distribution on the
product test space. Positivity follows from the Hilbert inner product,
the spectral condition from \(U\), covariance from (b), and locality
from (a). The adjoint companions give the adjoint relations; a densely
defined operator with a densely defined adjoint restriction is closable.
Cyclicity holds because the scalar subfamily already generates \(\HH\).
Vacuum uniqueness is part of the representation hypothesis. Thus the
fields satisfy the stated Wightman requirements. Property (c) supplies
the additional Yang--Mills identification, which is not a consequence
of those axioms alone. Property (d) excludes a theory consisting only of
invariant vectors.\end{proof}

There is no need for constants in \eqref{eq:wordbound} uniform in \(n\)
to establish the temperedness of each Wightman distribution in this
already constructed representation. This does not say that arbitrary
fixed-order Euclidean estimates suffice to perform the original
reconstruction. That earlier step requires its own growth condition.

\subsection{A concrete tensor-covariance test}

For actual approximants \(\Phi_{D,G}\), a sufficient calibration estimate
is, uniformly for Lorentz transformations \(\Lambda\) in compact sets,
\begin{multline}
 \big\|(1+H)^{-p}\big[
 U(\Lambda)\Phi_{D,G}(v,f)U(\Lambda)^{-1}\\
 -\Phi_{D,G}(\rho_D(\Lambda)v,f\circ\Lambda^{-1})
 \big](1+H)^{-p}\big\|\longrightarrow0. \label{eq:covdefect}
\end{multline}
The approximants must converge in a compatible energy-form topology.
On compact Lorentz sets the forward-cone spectral condition makes the
transformed energies comparable, so the conjugations respect
\(\DD_\infty\). Passing to the limit in \eqref{eq:covdefect} proves
covariance. An alternative is to prove the corresponding ordered
Euclidean covariance and then continue analytically. Euclidean rotations
are not being represented here by unitary operators on the physical
Hilbert space.

A positive 21-dimensional Gram matrix establishes independence, not
\eqref{eq:covdefect}. The Lorentz orbit may otherwise leave the proposed
finite field space or carry a different representation. An intrinsic
finite-dimensional completeness theorem with an equivariant injective
symbol map would be another route; its hypotheses are not supplied by
Gaussian rank alone.

\subsection{What an ultraviolet identification estimate must say}

At a fixed nonexceptional rescaled configuration, a typical required
statement, after the declared canonical scaling and logarithmic mixing,
is
\begin{equation}
 S_{\mathrm{rescaled},\lambda}
  -\sum_{k=0}^K g(\lambda)^k S_k^{\mathrm{YM}}
 =O\!\left(g(\lambda)^{K+1}|\log g(\lambda)|^{q_K}\right)
 \label{eq:uvmatch}
\end{equation}
in a specified distribution or energy topology, for every fixed requested
order \(K\). Exceptional configurations require their own local OPE
subtractions and remainders. Equation~\eqref{eq:uvmatch} is a target
form, not a result established by the present one-loop calculation.
The latter fixes a useful leading coupling coordinate but does not
identify every composite field or contact normalization.

The first genuinely nonabelian algebraic test beyond the 21 quadratic
directions is an oriented triple. A purely Gaussian scalar construction
preserves improper color-reflection symmetry and cannot generate it.
The fixed-\(q\) interaction calculation supplies a candidate first-order
coefficient. To construct a field, its coefficient must be identified
for the actual long prefix with a remainder smaller than that coefficient
in the common operator topology. Dividing by a small coefficient before
proving a relative bound is not justified.

\subsection{The next mathematical tasks, in dependency order}

\begin{enumerate}
 \item \textbf{Attach the modified intrinsic minima.}
 Prove Problem~\ref{prob:admission} on preparation-changed slots,
 including complete paths, crossing plaquettes, retained gauges and
 simultaneous contour margins. The unchanged-slot identity
 Lemma~\ref{lem:slots} is already available. Apply the local current
 and complex-margin lemmas only after their patch hypotheses hold.
 \item \textbf{Close the literal all-spectator history estimate.}
 Use the original birth dates, root masks, causal periods and complete
 source-dependent scalar ledger. Include the native \(R\)-errors and
 the full marked \(T\)-operators, with all first-mark cross terms.
 The output must dominate the original normalized connected response,
 rather than an isolated regular numerator.
 \item \textbf{Complete the scalar continuum application.}
 Verify one hierarchy with the required all-order growth, positivity,
 continuous rotations, domain estimates and a positive-energy witness.
 Fixed-coupling, volume-first diagonal reflection positivity is useful
 but must be matched to the chosen ultraviolet limit; a different
 thermodynamic state cannot be substituted without proof.
 \item \textbf{Construct and calibrate the quadratic frame.}
 Establish the actual relative 21-field Gram bound, the matching matrix
 and the convergent heat/energy limits, then prove tensor covariance
 by \eqref{eq:covdefect} or a genuine completeness theorem.
 \item \textbf{Close cubic transport and production.}
 Prove dated \(h^6\) transport of the rooted cubic projection, paying
 the \(h^{-4}\) root census, and control the normalized remainder
 feedback in one common spatial reserve. Match actual production to
 the nonzero first-order coefficient with a relative remainder.
 Include annular errors, winding tails and equation-of-motion contact
 terms. One-step constants alone do not yield this induction.
 \item \textbf{Extend to the higher observable hierarchy.}
 Add curvature jets and renormalized products with energy bounds,
 locality, tensor calibration, and the Yang--Mills matching
 \eqref{eq:uvmatch}. Apply Theorem~\ref{thm:completion} to the resulting
 compatible fields. This completes the stated existence goal if all
 its premises are verified; a mass gap remains a further theorem.
\end{enumerate}

The list specifies sufficient work, not a claim that every item is
logically necessary for every conceivable construction. A different
route may replace an item by a theorem with the same conclusion. It
must preserve the actual measure, the common state, the full observable
algebra and the normalization which the present route requires.

\researchfile{work/full_theory_after_21_obstruction_and_completion_2026-10-18.txt}
